\section{Introduction}

This paper stays within the same research line as the existing Fishery and Harvest work: renewable commons under strategic pressure. The goal is not to replace those environments with a new domain, but to make the institutional interpretation clearer and the bridge to real governance questions more credible.

The framing follows two linked observations. First, common-pool resource governance is not a single design problem. Different settings rely on different mixes of centralized intervention, local monitoring, and multi-level coordination \citep{ostrom1990, hilborn2005, nagendra2012}. Second, recent work on LLM populations shows that collective outcomes can deteriorate as strategy populations become more exploitative or competitively selected \citep{willis2025cooperation, willis2026collective}. The natural synthesis is to ask not only whether cooperation emerges, but which institutions remain robust when strategic pressure intensifies.

The contribution here is therefore a controlled institutional-design extension of Harvest Commons. Rather than claiming field calibration, the paper uses literature-backed scenario archetypes that correspond to recognizable governance forms in fisheries, irrigation, and forest co-management. These archetypes are paired with explicit oversight frictions such as imperfect detection, delayed enforcement, and limited intervention capacity. The resulting study is intended as a policy-interpretable benchmark for comparing institutional designs before domain-specific field study or deployment.
